% !TITLE 饮中八仙歌
% !AUTHOR 杜甫
% !DYNASTY 唐
% !ORDER 20

\begin{poem}
知章\textsuperscript{1}骑马似乘船，眼花落井水底眠。
汝阳\textsuperscript{2}三斗始朝天，道逢麴车\textsuperscript{3}口流涎，恨不移封向酒泉\textsuperscript{4}。
左相\textsuperscript{5}日兴费万钱，饮如长鲸吸百川，衔杯乐圣称避贤\textsuperscript{6}。
宗之\textsuperscript{7}潇洒美少年，举觞白眼望青天，皎如玉树临风前\textsuperscript{8}。
苏晋\textsuperscript{9}长斋绣佛前，醉中往往爱逃禅\textsuperscript{10}。
李白斗酒诗百篇，长安市上酒家眠，天子呼来不上船，自称臣是酒中仙。
张旭\textsuperscript{11}三杯草圣传，脱帽露顶王公前，挥毫落纸如云烟。
焦遂\textsuperscript{12}五斗方卓然，高谈雄辩惊四筵。
\end{poem}

\begin{annotations}
\notetext{1}{知章：贺知章（659-744），盛唐诗人，以狂放善饮著称。他喝醉了骑马摇晃如坐船，眼花掉进井里，竟然在井底睡着了。}
\notetext{2}{汝阳：汝阳王李琎，唐玄宗的侄子。喝了三斗酒才去朝见天子。}
\notetext{3}{麴车：运酒曲（酿酒的酵母）的车。路上碰到酒车就流口水。麴（qū）。}
\notetext{4}{酒泉：地名，在今甘肃。传说此地有泉味如酒。“恨不移封向酒泉”——恨不得把封地改到酒泉去。}
\notetext{5}{左相：李适之，曾任左丞相（唐朝最高官职之一）。每天喝酒花费万钱，饮酒量像长鲸吸干百条河流。}
\notetext{6}{衔杯乐圣称避贤：李适之被李林甫排挤罢相后，作诗自嘲“避贤初罢相，乐圣且衔杯”。乐圣，爱酒（古代称清酒为“圣人”）。}
\notetext{7}{宗之：崔宗之，名士崔日用之子。举杯时翻白眼望天（魏晋名士以白眼表示不屑流俗），风度像玉树在风前摇曳。}
\notetext{8}{玉树临风前：像玉树一样挺立在风前，形容仪态高洁优美。这个成语迄今还在使用。}
\notetext{9}{苏晋：进士出身，善属文。长期在家中的绣佛像前吃斋，但一喝酒就把佛门戒律忘光了。}
\notetext{10}{逃禅：逃离禅规、破戒。}
\notetext{11}{张旭：盛唐最著名的草书书法家，称“草圣”。三杯酒后，当着王公脱帽露顶（不拘礼法），挥毫如云烟。}
\notetext{12}{焦遂：平民，口吃，但喝了五斗酒后精神焕发，高谈阔论令人惊服。卓然，神采焕发的样子。}
\end{annotations}

\begin{appreciation}
《饮中八仙歌》是杜甫写得最开心的一首诗，也是中国文学史上最热闹的群像诗之一。杜甫一辈子写诗，多数时候沉着脸，只有这一首，从头到尾像在讲段子——八个酒友，一人两三句，写完一个笑一个。

贺知章骑马似乘船——喝多了，在马背上晃得像坐船；眼花落井水底眠——掉井里，居然在井底睡着了。汝阳王李琎，三斗酒下肚才去上朝，路上碰到运酒的车，口水直流，恨不移封向酒泉——恨不得把封地搬到酒泉去。左相李适之，喝酒像长鲸吸百川，鲸鱼喝水都没他量大。崔宗之是个美少年，举着酒杯翻白眼望天，皎如玉树临风前——玉树临风这个成语，现在还在用，出处就在这里。苏晋最可爱：白天在绣佛前吃斋，一喝醉就逃禅破戒。张旭三杯酒下肚，敢在王公面前脱帽露顶，写字如云烟——他是草圣。焦遂平时口吃，五斗酒下肚，高谈雄辩惊四筵。

八个人里，李白独占四句：李白斗酒诗百篇，长安市上酒家眠，天子呼来不上船，自称臣是酒中仙。天子呼来不上船——皇帝叫都不去。书里李白的《将进酒》说"人生得意须尽欢"，李白自己就是这句话的活人版。酒是李白的燃料。

还记得《鹿鸣》吗？那是中国最早的宴饮诗，朋友们聚在一起，鼓瑟吹笙。一千多年后，杜甫给盛唐的名士们补了一张酒桌——这八个人里，有王爷，有宰相，有诗仙，有草圣。盛世的长安，允许人狂，允许人有个性，连皇帝都拿醉鬼没办法。

这张酒桌后来散了。贺知章早早去世，安史之乱一来，李白流放夜郎，长安换了天地。杜甫把这些人写进诗里，等于替一个时代拍了张合影——照片里的人醉着，时代醒着。

你以后也会有一群这样的朋友——不是能喝八斗酒的那种，是你在他们面前可以不用端着的那种。遇见了，要珍惜。酒就别学了，你自己也知道为什么。
\end{appreciation}
